\documentclass[10pt,a4paper]{amsart}
\usepackage[margin=19mm]{geometry}
\usepackage{amsmath,amssymb,mathtools}
\usepackage[scr=rsfs,cal=euler]{mathalfa}
\usepackage{xfrac}
\usepackage[unicode,hidelinks,bookmarks=false]{hyperref}
\newcommand{\ie}{i.e.\ }
\newcommand{\cf}{cf.\ }
\newcommand{\resp}{respectively\ }
\newcommand{\Subsection}{\subsection}
\providecommand{\nobreakdash}{\nobreak\hbox{-}}
\setlength{\emergencystretch}{2em}
\multlinegap0pt
\usepackage{bm}
\usepackage{bbm}
\usepackage{dsfont}
\usepackage{xspace}
\usepackage{cancel}
\usepackage{mathtools}
\usepackage{amsthm}
\makeatletter
\@ifundefined{lemma}{\newtheorem{lemma}{Lemma}}{}
\@ifundefined{problem}{\newtheorem{problem}{Problem}}{}
\@ifundefined{theorem}{\newtheorem{theorem}{Theorem}}{}
\makeatother
\def\LO{L^2(\O)}
\def\O{\Omega}
\def\PiK{\Pi_{K}^{\Delta}}
\def\l{\lambda}
\title[A posteriori error estimates for a $C^1$ virtual element met]{A posteriori error estimates for a $C^1$ virtual element method applied to the thin plate vibration problem}
\author{Franco Dassi, Andres E. Rubiano, Iv\'an Vel\'asquez}
\date{Source: arXiv:2507.06846v1 (CC BY 4.0)}
\begin{document}
\maketitle

\noindent\emph{Source and licence.} A posteriori error estimates for a $C^1$ virtual element method applied to the thin plate vibration problem. Version arXiv:2507.06846v1 (CC BY 4.0), distributed under the Creative Commons Attribution 4.0 International licence, as recorded in the arXiv licence metadata for this exact version and shown on the abstract page https://arxiv.org/abs/2507.06846v1. Full rights notice: licenses/arxiv-2507.06846-CC-BY-4.0.txt.\par\smallskip

\noindent\textit{Editorial scope.} Counted unit: the Theorem whose source label is \texttt{Teor4.9} in the article (source0.tex lines 841--847, source label \texttt{Teor4.9}), delivered together with the article's own complete original proof of it (source0.tex lines 848--905, one \texttt{proof} environment of 58 lines). No mathematics is written, repaired or re-derived: statement and proof are the article's own text line for line. Required context retained uncounted: WMPr (lines 300--305); DiscVP (lines 555--560); double\_conv\_2 (lines 748--751); Lema4.12 (lines 819--823); aux\_pre\_th (lines 820--822). Every definition, display and label of those lines is retained unchanged. Omitted from this package and not redistributed: 1--299, 306--554, 561--747, 752--818, 823--840, 906--1442. Nothing inside a retained excerpt is omitted or altered: every retained body is delivered exactly as the article's lines are written, so no omission receipt of any kind accompanies this package. Citations: the retained excerpts carry no citation command at all, so no bibliography entry of the article is transcribed and no reference is redistributed. Reference adaptation: none was needed and none was made; every reference inside the delivered unit resolves inside it, so no unresolved reference or citation is produced. Printed slips kept literal: no printed word, symbol, hypothesis or number of the counted statement or of its proof was altered. Layout and class: the article's own class is \texttt{elsarticle} with packages including mathrsfs, amssymb, amsmath, amsfonts, latexsym, amsthm, graphicx, float, epsfig, color, fancyhdr, framed, todonotes, comment; the wrapper is the lane's pinned \texttt{amsart} editorial wrapper at 10pt on A4, which restates the theorem-like environments the retained text needs and adds the article's own macro definitions that the retained text needs (\texttt{\textbackslash LO}, \texttt{\textbackslash O}, \texttt{\textbackslash PiK}, \texttt{\textbackslash l}), with extra wrapper packages (bm, bbm, dsfont, xspace, cancel, mathtools). The delivered numbering is the unit's own. Assets: no retained span contains a figure environment, a graphics-inclusion command, a PDF-page inclusion, a table or a TikZ picture; no image file is copied into this package, the package declares no asset dependency and no image file is redistributed with it. Licence: version arXiv:2507.06846v1 of this article is distributed under the Creative Commons Attribution 4.0 International licence, as recorded in the arXiv licence metadata for this exact version and shown on the abstract page https://arxiv.org/abs/2507.06846v1. A full rights notice is delivered at licenses/arxiv-2507.06846-CC-BY-4.0.txt. Characters: every retained excerpt is pure ASCII, so the delivered engine, XeLaTeX with its default Latin Modern text font, reports no missing character.
\begin{problem}\label{WMPr} 
Find $(\lambda,u)\in \mathbb{R}\times V$, with $u\neq 0$ such that
\begin{equation*}
    a(u,v)=\lambda b(u,v)\qquad \forall v\in V.
\end{equation*}
\end{problem}

\begin{problem}\label{DiscVP}
Find $(\lambda_h,u_h)\in \mathbb{R}\times V_h$, with $u_h\neq 0$ such that
\begin{equation*}
    a(u_h,v_h)=\lambda_h b_h(u_h,v_h)\qquad \forall v_h\in V_h.
\end{equation*}
\end{problem} 

\begin{align}
\left|\l-\l_h^{(j)}\right|&\lesssim  h^{2s},\qquad j=1,\ldots,m\,.\label{double_conv_1}\\[1ex]
\left|\l-\l_h^{(j)}\right|&\lesssim \Big( ||u-u_h||_{2,\Omega}^2 + ||u-\Pi_K^0 u_h||_{0,\Omega}^2 + |u- \PiK u_h|_{2,h}^2 \Big)\label{double_conv_2}.
\end{align}

\begin{lemma}\label{Lema4.12} Let $(\mu_h^{(j)},u_h)$ be an eigenpair of $\widehat{T}_h$  with $j=1,2,...,m$ and $||u_h ||_{0,\O}=1$. Then, there exist an eigenfunction $u\in \LO$ of $T$ associated to $\mu$, and a positive constant $C$ such that
\begin{align}
||u-u_h ||_{0,\O}\leq  Ch^s\Big\{|w-w_h|_{2,\O}+||w-\Pi_K^0 w_h ||_{0,\Omega}+|w-\PiK  w_h|_{2,h}\Big\} \label{aux_pre_th}
\end{align}  
\end{lemma}

\begin{theorem}\label{Teor4.9}
There exist $s\in(1/2,1]$ and a positive constant $C$ independent of parameter $h$ such that
\begin{align*}
||u-u_h ||_{0,\O}\leq Ch^s\Big\{|u-u_h|_{2,\O}+||u-\Pi_K^0 u_h ||_{0,\Omega}+|u-\PiK  u_h|_{2,h}\Big\}.
\end{align*}
%with $s$ like in Lemma~\ref{lem:uZero}.
\end{theorem}

\begin{proof}
The proof proceeds by bounding each term on the right-hand side of  \eqref{aux_pre_th} in  Lemma~\ref{Lema4.12}. We first consider $w\in V$ and $w_h\in V_h$ as the solutions of problems
\begin{align}
a(w,v)=b(u,v), \quad
a_h(w_h,v_h)=b_h(\mathcal{P}_hu,v_h), \quad \forall v\in V,\, \forall v_h \in V_h \label{aux_disc_prob_apost},
\end{align}
respectively. Since $b(u,v)=\frac{1}{\lambda}a(u,v)$ (cf. Problem~\ref{WMPr}) we get $w=\frac{1}{\lambda}
u$ with $\lambda\neq 0$ . Thus, 
\begin{align}
|w-w_h|_{2,\O}&\leq \lambda^{-1}|u-u_h|_{2,\O}+|\lambda^{-1}-\lambda_h^{-1}||u_h|_{2,\O}+|\lambda^{-1}u_h-w_h|.\label{aux_1ap} 
\end{align}
Note that, from estimate~\ref{double_conv_2} we obtain
\begin{align}
|\lambda^{-1}-\lambda_h^{-1}|=|\lambda\lambda_{h}^{-1}||\lambda-\lambda_h|\leq C\Big( ||u-u_h||_{2,\Omega}^2 + ||u-\Pi_K^0 u_h||_{0,\Omega}^2 + |u- \PiK u_h|_{2,h}^2 \Big)\label{est_chu}.
\end{align}

On the other hand, we know that $a_h(\lambda^{-1}u_h,v_h)=b_h(u_h,v_h)\ \forall v_h \in V_h $ (cf. Problem~\ref{DiscVP}). Hence, subtracting this equality from \eqref{aux_disc_prob_apost} we obtain 
\begin{align*}
a_h(w_h-\lambda^{-1}u_h,v_h)=b_h(\mathcal{P}_hu-u_h,v_h) \ \forall v_h \in V_h .
\end{align*}
Therefore, using the fact that $a_h(\cdot,\cdot)$ is $V_h-$elliptic, we can deduce
\begin{align}
|w_h-\lambda^{-1}u_h|_{2,\O}^2&=a_h(w_h-\lambda^{-1}u_h,w_h-\lambda^{-1}u_h)=Cb_h(\mathcal{P}_hu-u_h,w_h-\lambda^{-1}u_h)\nonumber\\
&\leq C || \mathcal{P}_hu-u_h||_{0,\O}||w_h-\lambda^{-1}u_h ||_{0,\O}\leq C || \mathcal{P}_h(u-u_h)||_{0,\O}||w_h-\lambda^{-1}u_h ||_{0,\O}\nonumber\\
&\leq C ||u-u_h||_{0,\O} ||w_h-\lambda^{-1}u_h ||_{2,\O}\leq C ||u-u_h||_{2,\O} |w_h-\lambda^{-1}u_h |_{2,\O}\nonumber\\
&\leq C |u-u_h|_{2,\O} |w_h-\lambda^{-1}u_h |_{2,\O},\nonumber
\end{align}
which implies that
\begin{align}
|w_h-\lambda^{-1}u_h|_{2,\O}&\leq C |u-u_h|_{2,\O} .\label{cua_chu}
\end{align}
Then, replacing \eqref{est_chu} and \eqref{cua_chu} in the right-hand side of \eqref{aux_1ap} lead to
\begin{align}
|w-w_h|_{2,\O}\leq C\Big( ||u-u_h||_{2,\Omega} + ||u-\Pi_K^0 u_h||_{0,\Omega} + |u- \PiK u_h|_{2,h}\Big)\label{1stToBound}.
\end{align}

To bound the second term on the right-hand side of \eqref{aux_pre_th} we apply the triangle inequality and we have
\begin{align}
||w-\Pi_K^0 w_h||_{0,\Omega}&\leq |w-w_h|_{2,\O}+||w_h-\Pi_K^0  w_h||_{0,\Omega}\label{like_c.vs_t}.
\end{align}
Next, for the second term in \eqref{like_c.vs_t}, we obtain
\begin{align}
||w_h-\Pi_K^0  w_h||_{0,\Omega}&\leq ||w_h-\lambda_h^{-1}u_h||_{0,\O}+\lambda_h^{-1}||u_h-\Pi_K^0 u_h||_{0,\Omega}+||\Pi_K^0(\lambda_h^{-1}u_h-w_h)||_{0,\Omega}\nonumber\\
&\leq 2|w_h-\lambda_h^{-1}u_h|_{2,\O}+\lambda_h^{-1}||u-u_h||_{0,\O}+\lambda_h^{-1}||u-\Pi_K^0 u_h||_{0,\Omega}\nonumber\\
&\leq C\Big\{||u-u_h||_{0,\O}+||u-\Pi_K^0 u_h||_{0,\Omega}\Big\}\label{car_fel}.
\end{align}
Thus, replacing  \eqref{1stToBound} and \eqref{car_fel} in \eqref{like_c.vs_t} imply that
\begin{align}
||w-\Pi_K^0  w_h||_{0,\Omega}\leq C\Big\{|u-u_h|_{2,\O}+||u-\Pi_K^0  u_h||_{0,\Omega}+|u-\PiK  u_h|_{2,h} \Big\}\label{2ndTermToBound}.
\end{align}

The third term on the right-hand side of \eqref{aux_pre_th} is bounded by following the same arguments as those applied to obtain \eqref{2ndTermToBound}. Therefore,
\begin{align}
|w-\PiK w_h|_{2,h}&\leq |w-w_h|_{2,\Omega}+|w_h-\PiK   w_h|_{2,h}=|w-w_h|_{2,\Omega}+\sum\limits_{K\in \Omega_h}|w_h-\PiK   w_h|_{2,K}\nonumber\\
&\leq C\Big\{|u-u_h|_{2,\Omega}+ ||u-u_h||_{2,\O}+|u-\PiK u_h|_{2,h}\Big\}\label{3TermToBound}
\end{align}
Finally, substituting  \eqref{1stToBound}, \eqref{2ndTermToBound} and \eqref{3TermToBound} on the right-hand side of \eqref{aux_pre_th} concludes the proof.
\end{proof}

\end{document}
